\section{CONCLUSIONES}

Se desarrolló un tacómetro digital con control de velocidad sobre el LPC1769 que integra
los principales periféricos vistos en la materia: temporizadores en modo contador y con
interrupción por match, PWM, ADC en modo ráfaga con transferencia por DMA, UART por
interrupciones con buffers circulares y GPIO. El uso de DMA y del contador de TIMER1
permite que la adquisición de las señales analógicas y el conteo de pulsos se realicen sin
intervención de la CPU, que sólo atiende tres interrupciones breves y un lazo principal
sin esperas.

La organización del firmware en capas de aplicación, servicios y dispositivos resultó útil
durante la puesta en marcha: cambiar la UART0 por la UART3 y el canal PWM1.1 por el PWM1.2
por problemas en pines de la placa sólo requirió modificar los módulos de dispositivos,
sin tocar la lógica de la aplicación. Del mismo modo, concentrar los parámetros en
encabezados de configuración permite ajustar el umbral de alarma, las ganancias del
regulador o la tensión de referencia sin modificar el código.

Compilar los mismos fuentes para la PC con dispositivos simulados permitió probar la
lógica de modos, la alarma con histéresis, el regulador PI y el protocolo de comunicación
antes de disponer del hardware completo, y detectar errores que en la placa habrían sido
más difíciles de encontrar.

Entre las limitaciones del diseño se encuentran la resolución de 2\,RPM de la medición,
consecuencia de la ventana de 500\,ms y del disco de 60 ranuras, que podría mejorarse
midiendo el período entre pulsos con el modo captura del temporizador; la ausencia de
aislación galvánica entre control y potencia; y que las ganancias del regulador PI son
valores iniciales que deben ajustarse con ensayos sobre el motor real.

\pendiente{agregar los resultados de los ensayos en placa: respuesta del regulador ante
cambios de consigna, verificación de la alarma y observaciones de la puesta en marcha}.
